% ---------------------------------------------------------------------------
%  Preface
% ---------------------------------------------------------------------------
\chapter*{Preface}
\addcontentsline{toc}{chapter}{Preface}
\markboth{Preface}{Preface}

There is an old and stubborn problem in the first year of every computer science
and computer engineering degree. Two courses run in parallel, and they seem to
belong to two different universes. In the computer architecture course, students
meet registers, instructions, and the bare metal of a processor; the work is
precise, but the reward is often nothing more than a number printed on a
terminal. In the programming course, students meet variables, functions, and
the comfortable abstractions of a high-level language; the work feels powerful,
but the machine underneath has vanished from view. A student can pass both
courses and still not feel, in any concrete way, that the \texttt{addi}
instruction they traced by hand on Monday is the same machine that ran their C
program on Thursday.

This book was written to dissolve that separation, and it does so with a
deliberately old-fashioned trick: it makes a game move on a screen. When a
white rectangle slides left because a student set a bit in a register, the
distance between the instruction and the consequence collapses to nothing. The
screen is honest in a way a grade is not. If the paddle leaves the viewport, the
boundary check was wrong, and no amount of confident explanation will put it
back. Retro game development turns the abstract machinery of a processor into
something a beginner can see, hear, and argue about.

The platform that makes this possible is \prg, an open educational runtime for
\riscv assembly and C games \citep{prg32repo}. \prg is deliberately not a CPU
emulator that hides the processor behind a comfortable layer. Student code runs
\emph{natively} as 32-bit \riscv machine code, either on a real ESP32\nobreakdash-C6
board on the desk or on Espressif's QEMU build on a laptop. The same three
functions---\code{init}, \code{update}, and \code{draw}---structure every
example, whether it is written in hand-tuned assembly or in plain C. Because the
two language tracks call exactly the same runtime, a teacher can show the same
game logic at two altitudes: as a C function that reads almost like prose, and as
a sequence of \riscv instructions whose every register can be traced. That dual
view is the pedagogical heart of this book.

Two practical commitments run through every chapter. First, every example is
designed to work on \emph{both} QEMU and the physical ESP32\nobreakdash-C6 board, so a
student without hardware is never locked out and a student with hardware is never
limited to simulation. Second, the work is incremental and visible. The intended
rhythm, borrowed directly from the \prg learning path, is small and concrete:
change one thing, run it, observe the result, and write down what changed.
Nothing in this book asks a beginner to type a hundred lines of assembly before
seeing a single pixel.

A word on why the processor at the centre of this book matters beyond the
classroom. \riscv is an open instruction set architecture, free of the licensing
restrictions that govern the proprietary architectures that have dominated
computing for decades. For Europe in particular, that openness has become a
matter of strategic technological sovereignty, and the introduction that follows
makes the case for treating \riscv literacy as an investment, not a curiosity.
Teaching a first-year student to read and write \riscv is, in a small way,
participation in that larger project.

The book is organised so that the two courses it serves can use it independently
or together. The computer architecture track and the computer programming track
occupy separate parts and never depend on each other's internals, exactly as the
\prg teaching guide recommends \citep{prg32teaching}. A short shared foundation at
the start, and a short reunion at the end, are the only places the two paths
touch. The appendices collect the reference material---the full runtime API, the
performance-measurement workflow, complete worked examples, the use of GitHub,
the construction of the physical board, and the configuration of the development
environment on Windows, Linux, and macOS---so that the chapters can stay focused
on teaching.

I hope that by the last page a student will not merely have passed two courses,
but will have built something that moves, makes noise, and runs on a chip they
can hold in their hand---and will understand, instruction by instruction, why it
works.
